\section{Evaluation}
\label{sec:eval}
A document can appear in the catalogue yet be missing from the index. A
search can reach the right collection yet miss the passage that answers the
question. We test these faults separately, and we treat answer quality as a
question of its own because it also depends on how much text the answering
agent can read. Section~\ref{sec:exp-e1} asks whether changes to a stored
document reach the index (RQ1). Section~\ref{sec:exp-e2} compares answer
quality in the native configurations (RQ6). Sections~\ref{sec:exp-e3}, \ref{sec:exp-e4}
and~\ref{sec:exp-e6} examine the stages that an answer depends on: collection
selection (RQ2), evidence-pack assembly (RQ3) and verifier-score ordering
(RQ4). Section~\ref{sec:exp-e5} reports checks of the experience store, and
Section~\ref{sec:exp-e7} examines answer traces, abstention and operational
failures (RQ5).

\subsection{Datasets and measures}\label{sec:eval-setup}
The answer comparison uses ten developer-authored questions with annotated
reference papers in five arXiv collections. The routing and evidence-pack
experiments use 72 in-domain questions across ten collections (138 documents
in the frozen run); out-of-domain probes are analysed separately. The
verifier diagnostics use a separate industrial corpus of 42 documents in
five collections. None of the saved questions contains a gold document
identifier.

Routing and evidence packs use different notions of a hit. For routing, a
hit means that the selector chose the labelled collection. For evidence
packs, it means that the annotated target string occurs in the returned text.
The second measure does not verify a claim, and it can count a match in a
document other than the annotated source. Hit@5 records whether a gold
document appears among the first five results. Answer scores and verdicts
are defined in Section~\ref{sec:e10}.

The paired sign and McNemar tests are two-sided and exploratory. Repeated
runs use the same questions, so they do not add independent observations.

\subsection{Do changes in storage reach the index? (E1)}
\label{sec:e1}
\label{sec:e2}
\label{sec:exp-e1}
We first checked the state that search depends on. After explicit indexing
and reindexing, all 328 stored arXiv parts had content-derived descriptions,
and search probes returned identifiers for all 100 papers. Of the 100
probes, 72 lacked collection scope and 45 returned a chunk-copy name first.
This establishes identifier presence, not ranking within each collection. The controlled
checks passed for document creation, moves, deletion and forced reindexing.

Two observations show why a search result alone does not locate a fault. In
an earlier test on a disposable collection, a rename left 18 vector entries
instead of nine; forced reindexing restored the expected count, so the index
was at fault. In a separate, later retest, an old label appeared in the
search results although the index metadata was correct; the label came from
the retrieval cache.

Not every check succeeded. Task-status and update checks failed, and merge
and split were not tested. A historical tag audit also recorded orphan tags
and a cross-collection registry anomaly. Together, the checks distinguish
what is stored, what is indexed and what the search interface returns.
